\section{Microservice Implementation}
\label{sec:implementation}

The four services share one Maven build, one hexagonal layout and one module of contracts. Every state change is a
single conditional write to MongoDB, so a crash at any point leaves work that \code{control} can redo. This section
covers the shared parts, the indexer in depth as an example, and the other services in outline.

\subsection{Code structure and technologies}

The system is one multi-module Maven build in Java~25. It contains \code{common}, the four services, and two modules
that are never deployed: \code{bench}, with the load generator and result processing, and \code{checks}, with the API
checks. Every dependency and plugin version lives only in the parent \code{pom.xml} (Table~\ref{tab:stack}). Each
service builds one executable jar and runs on its own, outside Compose too.

\begin{table}[ht]
\centering
\small
\caption{Technologies and versions.}
\label{tab:stack}
\begin{tabularx}{\textwidth}{@{}L{2.3cm}Y@{}}
\toprule
Piece & Choice \\
\midrule
HTTP & Javalin 6.7, one port per service (7001 to 7004) \\
JSON & Jackson 2.22 with one shared \code{ObjectMapper}: snake\_case, ISO-8601 dates, unknown fields ignored. It is Javalin's default JSON mapper, which is why it replaces the Gson of the assignment's example \\
Storage & MongoDB 7.0 with the synchronous Java driver 5.13; the datalake on a Docker volume \\
Coordination & Hazelcast 5.7, one member and every service a client; ActiveMQ Classic 6.1.8 in mode \mode{activemq} \\
Packaging & One shaded jar per service and one \code{Dockerfile} on \code{eclipse-temurin:25-jre}; Docker Compose and nginx \\
Tests and benchmarks & JUnit 5.11, Testcontainers 2.0, JMH 1.37, XChart 3.8 for the benchmark charts \\
\bottomrule
\end{tabularx}
\end{table}

Each service keeps the hexagonal layout of Stage~1. \code{model} holds records and pure logic, \code{ports} the
interfaces that logic needs, \code{adapters} their MongoDB, broker, datalake and HTTP implementations, and
\code{commands} the use cases. \code{model} and \code{commands} import no Javalin, MongoDB or broker classes, so they
are tested with in-memory fakes. Each service's \code{Main} and factory wire the pieces through constructors, without
an injection framework. Every service opens its HTTP port first, so \code{/health} answers \code{503} while it waits,
and then retries MongoDB and the broker without limit. Startup order therefore does not matter.

\subsection{Shared contracts in \texorpdfstring{\code{common}}{common}}

Everything that two or more modules use has its literal name in \code{SPEC.md} and, in code, lives only in
\code{common} (Table~\ref{tab:common}). Two services therefore cannot disagree on a collection, a field, a state or a
message. \code{common} changes only in a pull request that all four members review. A write used by a single service,
such as claiming a book or publishing a batch, stays in that service's MongoDB adapter.

\begin{table}[ht]
\centering
\small
\caption{What \code{common} holds.}
\label{tab:common}
\begin{tabularx}{\textwidth}{@{}L{2.3cm}Y@{}}
\toprule
Area & Content \\
\midrule
Model & Book, batch and rebuild states, error and problem codes, and the two messages, \code{BookTask} and \code{BatchPublished} \\
Configuration & \code{TarantinoConfig} reads every \code{TARANTINO_*} variable with its default and checks its invariants, such as heartbeat $\leq$ lease / 3; an invalid value stops the service naming the variable \\
Terms & The Stage~1 \code{Tokenizer} and stopwords, used by the indexer to count terms and by \code{query} to read \code{q} and \code{author} \\
Datalake & \code{DatalakePaths}, with the \code{batch} layout and the write to \code{<final>.<uuid>.tmp} renamed with \code{ATOMIC_MOVE} \\
MongoDB & Collection and field names, the indexes, the writes ``Create the pointer'', ``Create the state'' and ``Complete a publication'', and the filter of the books that block a generation switch \\
Messaging & The \code{Broker}, \code{BookTasks} and \code{BatchAnnouncements} interfaces, with their Hazelcast and ActiveMQ implementations and clients that connect in the background and keep retrying \\
HTTP & Problem Details for every error, \code{/health}, and \code{BookIds.parse} for every \code{{book_id}} route \\
Operations & \code{TransitionLog}, one line per state change, from which the benchmarks measure; \code{CrashPoints} for the recovery tests; and \code{LeaseHeartbeat} \\
\bottomrule
\end{tabularx}
\end{table}

\subsection{Conditional writes and leases}
\label{sec:conditional-writes}

Every state change is one conditional write to a single MongoDB document. Its filter carries the state it expects and,
while a job holds the work, its \code{claim_id}. No service reads a document, decides in Java and writes it back. Two
replicas racing for the same book or batch can therefore never both win, and no transaction is needed.
\spec{5.1} lists all 23 writes. This is the one that publishes a batch:

\begin{lstlisting}
batches_g<w>.findOneAndUpdate(
  filter: {_id: B, status: "WRITING", claim_id: C,
           $expr: {$gt: ["$lease_until", "$$NOW"]}},
  update: [{$set: {status: "PUBLISHED", published_at: "$$NOW", book_ids: books}},
           {$unset: "lease_until"}],
  returnDocument: AFTER)
\end{lstlisting}

\begin{itemize}
\item \textbf{Server time.} Leases and expirations are compared with MongoDB's \code{$$NOW} inside update pipelines~\cite{mongodb},
  never with a process clock, so clock drift between containers or hosts cannot break a lease.
\item \textbf{One write, the right result.} When the next step depends on how the document ended up, the write is a
  \code{findOneAndUpdate} that returns it, never a write followed by a read. The published batch it returns becomes the
  \code{BatchPublished} message.
\item \textbf{The heartbeat.} A job that holds a lease starts a \code{LeaseHeartbeat}, which renews it every
  \code{TARANTINO_HEARTBEAT_SECONDS} (20\,s for a 60\,s lease) on its own virtual thread. Only a renewal that does not
  match loses the lease; one that throws, as when MongoDB does not answer for a moment, is retried on the next beat.
  Before each of its writes the job checks whether the lease is lost, and once it is lost it writes nothing else.
\item \textbf{Crashes.} If a process dies or MongoDB fails halfway, nothing is marked: the lease expires and
  \code{control} redoes the work. A job never retries a write that did not match; it leaves the work to whoever won.
\end{itemize}

\subsection{The indexer}
\label{sec:indexer}

The indexer turns downloaded books into metadata and postings and makes them searchable one batch at a time. Nothing a
batch writes is visible until a single write publishes it. Figure~\ref{fig:batch} shows one batch.

\textbf{Gathering a batch.} \code{BatchCollector} waits for a first message on \code{books-to-index} and then drains the
queue until it has \code{TARANTINO_INDEX_BATCH} books (100) or \code{TARANTINO_BATCH_TIMEOUT_SECONDS} (30\,s) have
passed since that first message. A book repeated in several messages is kept once, with its earliest
\code{enqueued_at}, so the benchmarks can measure its queue wait.

\textbf{Writing it.} The indexer reads the write generation, which is the target of a running rebuild or else the
current one, and ensures that generation's indexes before its first write to it. It opens the batch as \code{WRITING}
with a new \code{claim_id} and a lease, starts its heartbeat, and claims the books with one \code{updateMany} from
\code{DOWNLOADED} to \code{INDEXING}. Only the books that this write moved are indexed; another replica may have claimed
the rest. \code{HeaderParser} reads each header and the \code{Tokenizer} counts the terms of each body. A book whose
body is missing goes to \code{FAILED} with \code{NOT_IN_DATALAKE}. The indexer then writes \code{metadata_g<w>} with
one \code{bulkWrite} of upserts and \code{postings_g<w>} in \code{bulkWrite}s of 1{,}000 terms, publishes, marks the
books \code{PUBLISHED} and announces the batch.

\begin{figure}[ht]
\centering
\begin{tikzpicture}[
  font=\footnotesize,
  stepnode/.style={circle, draw=inkprimary, fill=white, minimum size=0.46cm, inner sep=0pt, font=\scriptsize\bfseries},
  steptext/.style={anchor=west, text width=9.0cm, inner sep=1pt, text=inkprimary},
  lease/.style={anchor=west, draw=modehazelcast, fill=modehazelcast!8, rounded corners=1pt, inner sep=2pt, font=\scriptsize, text=modehazelcast!80!black},
  crash/.style={anchor=west, draw=modeactivemq, fill=modeactivemq!8, rounded corners=1pt, inner sep=2pt, font=\scriptsize\ttfamily, text=modeactivemq!85!black},
]
\def\gap{0.6}
\draw[inkmuted, line width=0.8pt] (0,0) -- (0,{-8*\gap});
\foreach \i/\text in {
  1/{Drain \code{books-to-index}: up to 100 books or 30\,s},
  2/{Read the write generation; ensure its indexes},
  3/{Open the batch: \code{WRITING}, \code{claim_id}, lease, heartbeat},
  4/{Claim its books: \code{DOWNLOADED} to \code{INDEXING}},
  5/{Read and parse each book; a missing body fails the book},
  6/{Write \code{metadata_g<w>}: one \code{bulkWrite} of upserts},
  7/{Write \code{postings_g<w>}: \code{bulkWrite}s of 1{,}000 terms},
  8/{Publish: one conditional \code{findOneAndUpdate}},
  9/{Mark the books \code{PUBLISHED}; send \code{BatchPublished}}} {
  \node[stepnode] at (0,{-(\i-1)*\gap}) {\i};
  \node[steptext] at (0.4,{-(\i-1)*\gap}) {\text};
}
\foreach \i/\text in {5/{lease checked}, 6/{lease checked}, 7/{lease checked}, 8/{lease in the filter}} {
  \node[lease] at (9.55,{-(\i-1)*\gap}) {\text};
}
\node[crash] at (12.75,{-0.5*\gap}) {after-dequeue};
\node[crash] at (12.75,{-6*\gap}) {during-batch-write};
\node[crash] at (12.75,{-7.5*\gap}) {after-publish};
\node[crash] at (12.75,{-8.25*\gap}) {after-mark};
\node[anchor=south west, font=\scriptsize\bfseries, text=modehazelcast!80!black] at (9.55,{0.25}) {Lease checks};
\node[anchor=south west, font=\scriptsize\bfseries, text=modeactivemq!85!black] at (12.75,{0.25}) {Crash points};
\end{tikzpicture}
\caption{One batch in the indexer: nine steps, the lease checks (blue), made before each write of steps 5 to 7 and inside the
publication filter, and the four crash points of the recovery tests (red), placed where each one stops the service. \code{during-batch-write} stops after the first
postings \code{bulkWrite}, so the batch is really half-written, and \code{after-mark} stops before
\code{BatchPublished} is sent. If the lease is lost at any check, the batch stops and is never published; if no book is
left after step 5, the batch is left to expire. In both cases \code{control} abandons it and sends its books back to the
queue.}
\label{fig:batch}
\end{figure}

\textbf{One posting per term and batch.} Stage~1 kept one document per word and added books with \code{$addToSet}.
Here each batch writes its own document per term, with the batch's books in ascending order and each term frequency at
the same position. Three properties follow. A half-written batch is invisible, because \code{query} only reads
postings whose \code{batch_id} is in its snapshot of published batches. Replicas never write the same document, so
frequent words cause no contention and no document grows towards MongoDB's 16\,MB limit. And upserts make a repeated
write harmless. The cost moves to reading: a term has one document per batch that contains it, so the batch size trades
indexing throughput against the number of documents a search reads.

\textbf{Metadata.} \code{HeaderParser} keeps the title and author of Stage~1 and adds the fields of the search
filters. \code{year} is the year of the \code{Release date:} line, never the eBook number on the same line, which can
also have four digits. \code{original_year} comes from \code{Original publication:} when the header has it, which only
24 of the 2{,}800 corpus headers do. \code{languages} holds ISO~639-1 codes in header order, and \code{author_terms} the
author's words through the same \code{Tokenizer}.

\textbf{Rebuild.} Each generation of the index lives in its own \code{metadata_g<n>}, \code{postings_g<n>} and
\code{batches_g<n>}, and \code{index_state.current} says which one \code{query} serves. During a rebuild the indexer
writes into the target generation while \code{query} keeps serving the old one. \code{control} sends every book back to
be indexed and switches the pointer with one write once no book blocks the switch. Renaming three collections would
not be atomic; switching a pointer is.

\textbf{Known limitation.} \code{POST /index/rebuild} reads the current generation and then opens the rebuild, while
\code{control} switches the generation and then finishes the rebuild: two writes on each side. A request that falls
between \code{control}'s two writes answers \code{202} but builds no new generation, so the operator sees
\code{rebuild: null} in \code{GET /index/status} and sends it again. Closing the gap needs ``Open a rebuild'' to
compute its target inside its own write, a contract change left for Stage~3.

\subsection{Crawler, query and control}

The other three services follow the same pattern: a claim, a lease renewed by \code{LeaseHeartbeat}, conditional
writes, and a message only after the write.

\textbf{Crawler.} Each replica runs eight download workers on platform threads, and each worker takes one book at a
time from \code{books-to-download}. A worker claims the book (\code{QUEUED} to \code{DOWNLOADING}, with its
\code{claim_id}, a lease and one more attempt) and skips the download if the book's body is already in the datalake.
Otherwise it reads the text from the local mirror or over HTTP, splits it into header and body at the Gutenberg markers
and writes each file to a unique \code{.tmp}, renamed with \code{ATOMIC_MOVE}, header first. ``Finish a download'' then
marks it \code{DOWNLOADED}, and only after that write is the book enqueued in \code{books-to-index}. A failed download
goes back to \code{QUEUED} while attempts are left if the error may pass (\code{NETWORK_ERROR},
\code{STORAGE_ERROR}), and straight to \code{FAILED} if it cannot (\code{NOT_FOUND}, \code{MISSING_MARKERS}). When
Gutenberg's mirror answers \code{429} or \code{503}, every download of the replica waits what \code{Retry-After} asks,
or 1, 2, 4, 8 and 16\,s.

\textbf{Query.} \code{query} keeps an immutable snapshot of the published batches of the current generation behind an
\code{AtomicReference}. One thread refreshes it at startup, on each \code{BatchPublished} and every 5\,s, comparing
batches by \code{_id}, never by publication time. A search reads the snapshot once and tokenizes \code{q}. With
filters, it first asks \code{metadata_g<n>} for the books matching \code{author_terms}, \code{languages} and
\code{year}. It then reads only the postings of the snapshot's batches, intersects the terms starting with the rarest,
scores each book with the sum of $(1 + \ln tf) \cdot \ln(N/df)$, breaks ties by \code{book_id} and reads the metadata of
the page in one query. Every request therefore sees one consistent published state, and no lock is taken on the read
path.

\textbf{Control.} At startup \code{control} creates the state of every candidate and enqueues only the books it
created, so a restart never seeds twice. Then it reconciles every \code{TARANTINO_RECONCILE_SECONDS}, in six idempotent
steps: (1) release expired downloads; (2) abandon the \code{WRITING} batches whose lease expired, in every generation;
(3) settle each \code{INDEXING} book by its batch, releasing it if the batch is abandoned and completing its publication
if the batch is published; (4) requeue the books whose message was lost, when their queue is empty, which also refills
the queues after a broker restart; (5) delete the postings of batches abandoned more than \code{TARANTINO_GC_SECONDS}
ago, on every pass, so a zombie indexer's late writes are deleted too; and (6) drive a rebuild and drop old
generations. A separate task deletes orphan \code{.tmp} files, and \code{GET /control/status} reports the generation,
the books per state, both queues and the last reconciliation.
